\documentclass[11pt,a4paper]{letter}
\usepackage[margin=1in]{geometry}
\usepackage{hyperref}
\signature{Ittipon Fongkaew \\ Corresponding author}
\address{School of Physics, Institute of Science \\ Suranaree University of Technology \\ 111 University Avenue \\ Nakhon Ratchasima 30000, Thailand \\ \texttt{ittipon@sut.ac.th}}
\begin{document}

\begin{letter}{Editor-in-Chief \\ \emph{Applied Energy} \\ Elsevier}

\opening{Dear Editor,}

We are pleased to submit the manuscript \emph{``A Probabilistic 48-Step Forecasting Benchmark for the Japan Electric Power Exchange: Conformal, Distributional, and Quantile-Averaging Methods Across Day-Ahead and Imbalance Markets''} for consideration as a research article in \emph{Applied Energy}.

The work presents the first systematic probabilistic electricity-price-forecasting (EPF) benchmark on the Japan Electric Power Exchange (JEPX) that meets the methodological standard of the Lago, Marcjasz, De Schutter \& Weron (2021) open-access European benchmark published in this journal. We use a 365-day sliding training window with weekly recalibration to evaluate seven forecasters --- weekly-naive, LEAR, DDNN-Normal, DDNN-Johnson-SU, N-HiTS, QRA-rolling, and LEAR-plus-adaptive-conformal --- across two areas (Tokyo, Kansai) and two markets (day-ahead, imbalance) for a 305-day held-out test period (2023-03 to 2023-12, 14{,}640 issuance times \(\times\) 48 horizons per area-market pair).

The manuscript makes four contributions that we believe are of interest to the \emph{Applied Energy} readership:
\begin{enumerate}
\item \textbf{The first probabilistic Lago-protocol-compliant EPF benchmark on JEPX.} We close the methodological gap between the European EPF benchmark literature and Asian-market practice by porting the Lago 2021 rolling-window protocol to JEPX 30-min half-hourly data, extending it to a 48-step multi-output target, and adding distributional, ensembled, and conformally-calibrated probabilistic forecasters from the post-2020 EPF research frontier.
\item \textbf{A cross-market reversal of the optimal single model.} The Hansen-Lunde-Nason \emph{Model Confidence Set} at \(\alpha = 0.10\) contains exactly one model in each of the four (area, market) pairs: N-HiTS on both DA markets, LEAR on both IM markets. To our knowledge this is the first published instance of such a clean cross-market identification in EPF. The pattern is robust across areas (Tokyo and Kansai) and is directly actionable for JEPX market participants.
\item \textbf{A QRA-rolling ensemble that dominates on every metric.} Combining LEAR and DDNN-Normal as base learners with a 90-day rolling calibration and weekly refit, the QRA-rolling ensemble is the best point forecaster (rRMSE 0.75--0.84 across all four pairs) and the best probabilistic forecaster (CRPS 1.23--1.95; empirical coverage within 1 percentage point of nominal 90\% on three of four pairs). The 90-day rolling-calibration choice is itself a methodological recommendation: a single fixed 30-day calibration window under-covers the 90\% interval by 11--20 percentage points.
\item \textbf{An honest finding that the seasonal-naive baseline is exceptionally strong on JEPX-Tokyo DA.} Even rigorously rolling LEAR is 14.6\% worse than the weekly-naive baseline on this market --- an unusual pattern with respect to the European EPF literature. We attribute it to the unusual fidelity of the Tokyo duck-curve weekly repetition driven by industrial demand combined with high photovoltaic penetration, and we expect this finding to motivate market-structure work on JEPX.
\end{enumerate}

The earlier preprint of this work \citep*[on which the present manuscript is based but which it supersedes;][]{sriboriboon_v1e} contained a non-reproducible exponential-growth trading back-test that, on subsequent investigation by the present authors, traced to a look-ahead bias. The present manuscript explicitly retracts that result and replaces it with the transparent rolling-protocol evaluation described above. We mention this here in the interest of full disclosure; the v1e preprint is cited in the manuscript text and reference list.

To our knowledge no prior peer-reviewed study has provided this combination of (i) a JEPX-specific multi-area benchmark, (ii) Lago-protocol-compliant rolling-window recalibration, (iii) probabilistic forecasts with full CRPS / Winkler / coverage / MCS evaluation, and (iv) a public reproducibility package. We believe these elements together provide actionable methodological guidance to both EPF researchers and Japanese-market practitioners.

The work has not been published previously and is not under consideration elsewhere. All authors have approved the submission. We declare no competing interests.

We suggest the following potential reviewers, none of whom have collaborated with us in the past three years:
\begin{itemize}
\item Prof.\ Rafa\l{} Weron, Wroc\l{}aw University of Science and Technology (\texttt{rafal.weron@pwr.edu.pl}) --- author of the seminal EPF reviews and the Lago 2021 benchmark.
\item Dr.\ Jesus Lago, KU Leuven (\texttt{jesus.lago@kuleuven.be}) --- lead author of the Lago 2021 benchmark whose protocol we adopt.
\item Dr.\ Grzegorz Marcjasz, Wroc\l{}aw University of Science and Technology --- co-author of the DDNN paper that forms a methodological backbone of this work.
\item Prof.\ Tao Hong, University of North Carolina at Charlotte (\texttt{hongtao01@uncc.edu}) --- editor-in-chief of the IJF GEFCom competitions, broad EPF expertise.
\end{itemize}

Thank you for considering our submission. We look forward to your response and welcome any questions about the technical methodology or the reproducibility package.

\closing{Sincerely,}

\end{letter}
\end{document}
